% =====================================================================
% Figure contract
% 核心结论：学习从 grounding/action 先验进入方法谱系，并由 Data flywheel
%           把 rollout、验证、筛选、更新与新任务生成闭合为持续优化回路。
% 图原型：mode C / 三段横向结构；左侧先验链、中间方法桥、右侧五节点环形 hero。
% 证据层级：hero = 五节点顺时针闭环；支撑 = 左侧三卡与中桥；收束 = 底部注释条。
% 能不能更少？：不能；删卡会破坏先验链，删桥会丢方法谱系，删环节点会破坏闭环。
%
% Step ① 设计文档（Form B / Narrative，中等复杂度综述结构图）
% 0 参考：沿用 cua-survey-fig1-overview.tex 的低饱和配色、浅底 zone 与标题层级；
%   环形连线借鉴 feedback-loop.tex 的独立路径思想，不复刻其单条 U 型回路。
% 1 领域：Computer-Use Agents 学习与优化；中文叙述，技术术语保留英文。
% 2 格式：TikZ standalone；XeLaTeX + fontspec + xeCJK；全图 sans serif。
% 3 画布：目标印刷宽度 W_print = 21.0 cm；墨迹宽约 21.0 cm，外加 8 pt border
%   后总宽仍小于 22 cm；正文目标最小字号 7.5 pt。中桥因用户要求 122 字符单行，
%   采用窄体横向压缩并在 300 dpi 单独核验可读性，不改变其逐字内容。
% 4 模块：左 zone 4.2×5.2 cm，内含三张 3.55×0.74 cm 卡；中桥 9.7×0.96 cm；
%   右 zone 6.1×5.2 cm，五节点以 72° 等角分布；底栏 21.0×1.05 cm。
% 5 连线：左链两条短直箭头；policy→bridge 与 bridge→右 zone 两条短直箭头；
%   右环五条独立 bend-left 弧线，均按 rollout→验证→筛选→更新→新任务生成→rollout。
% 6 空间：三个横向区域共享中心轴；桥两端各留 0.5 cm 箭头通道；环节点外侧走弧线。
% 7 强调：右环用 teal 线与五种 pastel 节点色构成唯一 hero；桥和底栏用浅灰降噪。
% 8 密度：13 个可见结构元素，属于中等档；输入无数值，不嵌入 chart/heatmap。
% 9 Narrative：左 x∈[-10.5,-6.3] 为竖向先验链；中 x∈[-5.8,3.9] 为单行方法桥；
%   右 x∈[4.4,10.5] 为等角闭环；上方标题、下方注释条填满横向空白并收束语义。
% 10 可视化决策：以闭环拓扑本身承担视觉记忆点；不添加用户未指定的数值或说明。
%
% Pre-flight：P1 设计文档已写；P2 卡片用 positioning 构造、环用等角极坐标；
% P3 所有边的 source/destination 已列明；P4 N/A（无数据 viz）；P5 桥端通道独立；
% P6 zone 固定边界含足够 padding；P7 已对照 lessons、global rules 与 hard constraints。
% 拓扑记录：五节点闭环必须由五条独立弧线组成；每段 tip 均在下一节点一侧，
% 最后一段从“新任务生成”回到“rollout”，不使用单条 rounded-corners 长回路。
% =====================================================================

\documentclass[border=8pt]{standalone}
\usepackage{fontspec}
\usepackage{xeCJK}
\usepackage{xcolor}
\usepackage{graphicx}
\usepackage{tikz}
\setCJKmainfont{Heiti SC}
\setCJKsansfont{Heiti SC}
\newfontfamily\bridgefont{Arial Narrow}
\usetikzlibrary{arrows.meta,bending,positioning,calc,backgrounds,fit}

% Low-saturation academic palette, aligned with cua-survey-fig1-overview.tex.
\definecolor{ink}{HTML}{27364A}
\definecolor{muted}{HTML}{748092}
\definecolor{mainLine}{HTML}{4F667C}
\definecolor{mainFill}{HTML}{EAF0F4}
\definecolor{heroFill}{HTML}{F1F3F6}
\definecolor{heroBorder}{HTML}{3F566C}
\definecolor{blueLine}{HTML}{6D879E}
\definecolor{blueFill}{HTML}{E7EEF4}
\definecolor{purpleLine}{HTML}{817594}
\definecolor{purpleFill}{HTML}{EEEAF2}
\definecolor{goldLine}{HTML}{9B875B}
\definecolor{goldFill}{HTML}{F2EEE3}
\definecolor{tealLine}{HTML}{648D8A}
\definecolor{tealFill}{HTML}{E5F0EE}
\definecolor{roseLine}{HTML}{9A777B}
\definecolor{roseFill}{HTML}{F3E9EA}
\definecolor{greenLine}{HTML}{71896E}
\definecolor{greenFill}{HTML}{EAF0E8}

\tikzset{
  every node/.style={font=\sffamily\fontsize{8.2}{9.8}\selectfont,text=ink},
  zone/.style={
    rectangle,
    rounded corners=6pt,
    draw=heroBorder!72,
    fill=heroFill,
    line width=0.85pt,
    inner sep=0pt
  },
  zone title/.style={
    font=\sffamily\fontsize{9.4}{11.2}\selectfont\bfseries,
    text=heroBorder,
    anchor=north
  },
  prior card/.style={
    rectangle,
    rounded corners=4pt,
    minimum width=3.55cm,
    minimum height=0.74cm,
    align=center,
    inner sep=4pt,
    line width=0.65pt,
    font=\sffamily\fontsize{8.0}{9.4}\selectfont
  },
  bridge bar/.style={
    rectangle,
    rounded corners=4pt,
    draw=muted!52,
    fill=mainFill!72,
    line width=0.55pt,
    minimum width=9.70cm,
    minimum height=0.96cm,
    inner sep=4pt,
    align=center
  },
  fly node/.style={
    rectangle,
    rounded corners=4pt,
    minimum width=1.30cm,
    minimum height=0.70cm,
    align=center,
    inner sep=3pt,
    line width=0.72pt,
    font=\sffamily\fontsize{8.0}{9.4}\selectfont\bfseries
  },
  arrow short/.style={
    -{Stealth[length=3pt 1.0,width'=0pt 0.5,sep=0pt 0.3]},
    line width=0.8pt,
    shorten >=1pt,
    shorten <=0pt,
    color=mainLine
  },
  fly arrow/.style={
    -{Stealth[length=5pt 1.2,width'=0pt 0.58,sep=0pt 0.4]},
    line width=0.95pt,
    shorten >=0pt,
    shorten <=0pt,
    color=tealLine
  },
  summary bar/.style={
    rectangle,
    rounded corners=4pt,
    draw=muted!48,
    fill=mainFill!58,
    line width=0.55pt,
    minimum width=21.0cm,
    minimum height=1.05cm,
    inner sep=5pt,
    align=center
  }
}

\begin{document}
\sffamily
\begin{tikzpicture}[node distance=0.42cm]

% Two fixed-size zones; their combined width plus the bridge corridor is 21.0 cm.
\node[zone,minimum width=4.20cm,minimum height=5.20cm] (priorZone) at (0,0) {};
\node[zone,minimum width=6.10cm,minimum height=5.20cm,
      right=10.70cm of priorZone.east,anchor=west] (flyZone) {};

% Global title centered over the complete 21 cm span.
\coordinate (topCenter) at ($(priorZone.north west)!0.5!(flyZone.north east)$);
\node[
  anchor=south,
  font=\sffamily\fontsize{12.0}{14.2}\selectfont\bfseries,
  text=heroBorder
] at ([yshift=0.52cm]topCenter)
  {学习与优化：先验注入 → 方法谱系 → Data flywheel（§7）};

% Left panel: three cards linked vertically by short canonical arrows.
\node[zone title] at ([yshift=-0.25cm]priorZone.north) {先验注入 §7.1–7.2};
\node[prior card,fill=blueFill,draw=blueLine,
      anchor=north] (grounding) at ([yshift=-1.02cm]priorZone.north)
  {grounding pre-training};
\node[prior card,fill=purpleFill,draw=purpleLine,
      below=of grounding] (actionFT)
  {action fine-tuning};
\node[prior card,fill=goldFill,draw=goldLine,
      below=of actionFT] (initialPolicy)
  {初始 policy};
\draw[arrow short] (grounding.south) -- (actionFT.north);
\draw[arrow short] (actionFT.south) -- (initialPolicy.north);

% Middle bridge: exact user-supplied text kept on one physical line.
\node[bridge bar,anchor=west] (bridge)
  at ([xshift=0.55cm]priorZone.east |- initialPolicy.center)
  {\resizebox{9.25cm}{!}{\bridgefont\strut BC §7.3 · 偏好优化 §7.5 · Reward/Process model §7.6 · offline RL §7.7 · online RL §7.8 · RLVR §7.9 · rejection sampling §7.10}};
\draw[arrow short] (initialPolicy.east) -- (bridge.west);
\draw[arrow short] (bridge.east) -- (flyZone.west |- bridge.center);

% Right panel: five equal-angle nodes around a shared center.
\node[zone title] at ([yshift=-0.25cm]flyZone.north) {Data flywheel §7.11};
\coordinate (ringCenter) at ([yshift=-0.20cm]flyZone.center);
\node[fly node,fill=blueFill,draw=blueLine]
  (rollout) at ($(ringCenter)+(90:1.50cm)$) {rollout};
\node[fly node,fill=tealFill,draw=tealLine]
  (verify) at ($(ringCenter)+(18:1.50cm)$) {验证};
\node[fly node,fill=greenFill,draw=greenLine]
  (filter) at ($(ringCenter)+(-54:1.50cm)$) {筛选};
\node[fly node,fill=goldFill,draw=goldLine]
  (update) at ($(ringCenter)+(-126:1.50cm)$) {更新};
\node[fly node,fill=roseFill,draw=roseLine,minimum width=1.80cm]
  (newTask) at ($(ringCenter)+(162:1.50cm)$) {新任务生成};

% Five independent arcs: every tip follows the same clockwise direction.
\draw[fly arrow] (rollout) to[bend left=20] (verify);
\draw[fly arrow] (verify) to[bend left=20] (filter);
\draw[fly arrow] (filter) to[bend left=20] (update);
\draw[fly arrow] (update) to[bend left=20] (newTask);
\draw[fly arrow] (newTask) to[bend left=20] (rollout);

% Full-width annotation strip.
\node[summary bar,anchor=north] (summary)
  at ([yshift=-0.42cm]$(priorZone.south west)!0.5!(flyZone.south east)$)
  {\resizebox{20.25cm}{!}{\strut flywheel 的上限由 task validity、rollout 吞吐与 verifier 偏差决定（§1.3 阶段四）；推理期 planning·reflection·search 见 §7.14}};

\end{tikzpicture}
\end{document}
